\lyseis{Β1.5}{Μέτρηση γωνιών, σύγκριση και διχοτόμος}
\sol{1} Ελεύθερη κατασκευή. Η διχοτόμος των $40^\circ$, $75^\circ$, $130^\circ$ σχηματίζει $20^\circ$, $37{,}5^\circ$, $65^\circ$ με κάθε πλευρά.
\sol{2} α) $25{,}5^\circ$ \quad β) $48^\circ15'$ \quad γ) $3^\circ30'$ \quad δ) $90\cdot60'=5\,400'$.
\sol{3} $0{,}7^\circ<4{,}5^\circ<38^\circ<38^\circ30'<42^\circ$.
\sol{4} α) $35^\circ$ \quad β) $62^\circ$ \quad γ) $4^\circ15'$.
\sol{5} $56^\circ$.
\sol{6} α) $90^\circ$ \quad β) $180^\circ$ \quad γ) $90^\circ$ (η μεγάλη είναι $270^\circ$) \quad δ) $60^\circ$ \quad ε) $150^\circ$ \quad στ) $0^\circ$.
\sol{7} α) Σωστό. β) Λάθος. γ) Σωστό. δ) Σωστό. ε) Σωστό.
\sol{8} $84^\circ:2=42^\circ$.
\sol{9} Στις $5{:}00$ ($150^\circ$) και στις $7{:}00$ (γωνία $7\cdot30^\circ=210^\circ$, μικρότερη $360^\circ-210^\circ=150^\circ$). Υπάρχουν δύο ώρες.
\sol{10} $360^\circ:8=45^\circ$ και $360^\circ:6=60^\circ$.
\sol{11} $360^\circ:3=120^\circ$.
\sol{12} Διάβασε τη λάθος κλίμακα του μοιρογνωμονίου. Το σωστό είναι $180^\circ-140^\circ=40^\circ$.
\sol{13} Ο $360$ έχει πολλούς διαιρέτες ($2,\,3,\,4,\,5,\,6,\,8,\,9,\,10,\,12,\,\dots$), άρα ο κύκλος χωρίζεται εύκολα σε ίσα μέρη.
\sol{14} Όχι, γιατί το μέτρο εξαρτάται από το «άνοιγμα» και όχι από το μήκος των πλευρών.
\sol{15} Με το μοιρογνωμόνιο ή τη δίπλωση υπάρχει μόνο μία ημιευθεία στο εσωτερικό που σχηματίζει με τις πλευρές ίσες γωνίες.
\sol{16} $45^\circ$.
\sol{17} Ελεύθερη απάντηση.
\sol{18} Ελεύθερη απάντηση.
